\documentclass[11pt]{article}
\usepackage[margin=1in]{geometry}
\usepackage[T1]{fontenc}
\usepackage{lmodern}
\usepackage{microtype}
\usepackage{amsmath,amssymb}
\usepackage{booktabs}
\usepackage{graphicx}
\usepackage{xcolor}
\usepackage[numbers,sort&compress]{natbib}
\usepackage{authblk}
\usepackage{enumitem}
\definecolor{linkblue}{HTML}{2A78D6}
\usepackage[colorlinks=true,linkcolor=linkblue,citecolor=linkblue,urlcolor=linkblue]{hyperref}
\setlist{itemsep=2pt,topsep=4pt}

\newcommand{\nEpisodes}{1320}
\newcommand{\nLLMEpisodes}{0}
\newcommand{\nModels}{5}
\newcommand{\fitStatus}{The fitted simulator passed the held-out check (40 held-out participants vs 480 simulated; criterion: |d| < 0.5 and KS p > 0.01 for every summary).}
\newcommand{\resultsStatus}{No language-model episodes have completed yet; this section fills in automatically as the experiment engine runs. Baseline results are reported below.}

\newcommand{\fig}[3]{%
  \IfFileExists{figures/#1.pdf}{%
    \begin{figure}[t]\centering\includegraphics[width=#2\linewidth]{figures/#1.pdf}\caption{#3}\label{fig:#1}\end{figure}}{}}

\title{\textbf{Axiom: Do Language Models Design and Analyse\\Experiments Like Scientists?}\\[4pt]
\large A cognitive-control lab fitted to real human EEG, and a signalling lab on real single cells}
\author[1]{Sumeet Musfirah}
\affil[1]{Independent Researcher}
\date{}

\begin{document}
\maketitle

\begin{abstract}
Evaluations of AI scientists usually grade prose, often with another language model as the judge, and those judges can disagree at chance on the same outputs.
We instead build a laboratory whose answers are known.
Axiom-Neuro simulates a human cognitive-control study (a flanker task with working-memory load, transcranial stimulation, and EEG).
Each task hides a causal mechanism among six candidate edges, and an agent must find it by designing experiments under a fixed participant-hour budget, analysing the data with standard tools, and submitting calibrated claims.
Several tasks open with pilot data containing a planted confound that real human-subjects researchers fall for: a stimulation effect that is really practice, and a theta effect that is really blinks.
The simulator's noise, heterogeneity and artifacts are fitted to the ERP CORE flanker dataset (40 participants) and validated on held-out participants.
A second, semi-synthetic track uses real ERP CORE trials as the substrate and injects known effects on top.
Scoring needs no judge: claims are compared with the hidden mechanism using a proper scoring rule.
Because the same ideal analyst can be applied to any agent's data, each agent's error splits cleanly into an experiment-design part and an analysis part.
A second lab, Axiom-Bio, asks the same question on 7{,}466 real T cells \citep{sachs2005causal}: which signalling molecules directly regulate which, given real drug interventions; comparing real molecule names with abstract labels separates recalling textbook biology from inferring it from data.
Hypotheses, the stopping rule and all analyses were preregistered before any language-model episode ran, and the whole study runs unattended on free infrastructure with a public results ledger.
\resultsStatus
\end{abstract}

\section{Introduction}
Language models are increasingly asked to act as scientists: to propose experiments, run analyses and report conclusions \citep{jansen2024discoveryworld,gandhi2025boxinggym,chen2025scienceagentbench}.
The hard part of science is rarely the arithmetic.
It is judgement: choosing an experiment that can tell hypotheses apart, noticing that a factor is tied to session order, checking whether an EEG effect survives artifact removal, and stating conclusions with the confidence the data warrant.

These are exactly the skills that are hardest to grade.
A model's report can sound right and be wrong, and when another model is the grader, the grade inherits that model's blind spots.
We take the opposite approach: build a world where the truth is known, so that scoring is arithmetic.

\paragraph{Contributions.}
\begin{enumerate}
  \item \textbf{A judge-free scientific-reasoning environment at the human-experiment level.}
  Prior neuroscience discovery benchmarks work at the circuit or cell level \citep{sinha2026digitalfly,lueckmann2026zebrafish,murphy2026mda}.
  Axiom-Neuro targets population-level causal claims from noisy human participants, where factorial design, counterbalancing and artifact handling decide the answer.
  \item \textbf{Planted, ecologically real traps with mirror controls.}
  Each confound task (practice disguised as stimulation; blinks disguised as theta) has a held-out mirror task with the same pilot data and the opposite truth, so an agent cannot score by always distrusting the pilot.
  \item \textbf{Grounding in real data, two ways.}
  Nuisance parameters are fitted to ERP CORE \citep{kappenman2021erpcore} and validated on held-out participants; a semi-synthetic track replaces the parametric noise with real trials. Agreement between tracks is itself a preregistered test (H5).
  \item \textbf{A design-versus-analysis decomposition.}
  Re-analysing each agent's own data with an ideal analyst separates ``ran the wrong experiment'' from ``drew the wrong conclusion from the right one''.
  \item \textbf{A preregistered, self-running study.}
  Gates in code block language-model runs until the fit validates, the reference agents order correctly and the task set is frozen; an adaptive stopping rule decides how many replications each cell needs. Everything runs on free GitHub-hosted runners with local open-weight models.
\end{enumerate}

\section{Related work}
Simulated discovery with hidden ground truth is an active area.
BoxingGym \citep{gandhi2025boxinggym} offers ten generic environments scored by expected information gain and is the closest framework to ours.
DiscoveryWorld \citep{jansen2024discoveryworld} covers 120 text-world discovery tasks across eight topics.
ScienceAgentBench \citep{chen2025scienceagentbench} evaluates data-analysis code from 44 papers but not experiment design.
In neuroscience, recent work evaluates agents on a whole-brain \emph{Drosophila} model \citep{sinha2026digitalfly}, identifies mechanistic models in a simulated zebrafish \citep{lueckmann2026zebrafish}, couples language models with Bayesian experiment design \citep{murphy2026mda}, and stress-tests coding agents on a real optogenetics pipeline, finding that stage-level successes do not compose into end-to-end discovery \citep{horstmann2026casestudy}.
CheeseBench \citep{bugaud2026cheesebench} places the model in the role of the experimental animal; we place it in the role of the experimenter.
Our niche is the human cognitive-experiment level, with its characteristic confounds.

\section{Environment}
\paragraph{Hypothesis family.}
Six candidate causal edges among load $L$, stimulation $S$, practice $P$, performance $Y$, frontal theta $T$ and blinks $B$:
$L\!\to\!Y$, $S\!\to\!Y$ (a main effect), $S{\times}L\!\to\!Y$ (an interaction), $P\!\to\!Y$, $L\!\to\!T$ (a genuine neural effect) and $L\!\to\!B$.
Six task types (four training, two held out) fix which edges exist; effect magnitudes are sampled per instance, and practice is present at random except where a trap requires it.
Practice follows a saturating learning curve, so a linear time covariate cannot remove an order confound; only counterbalancing or a between-subjects design can.

\paragraph{Actions and budget.}
An agent may design an experiment (factors, fixed levels, within or between subjects, fixed or counterbalanced order, sample size, trials per cell, EEG and ICA), describe data, fit cluster-robust regressions, check order confounds, check artifacts, and submit a claim and confidence for every edge.
Experiments cost participant-hours (session length times participants, more with EEG); analysis is free.
The budget (30 participant-hours) affords a well-chosen factorial study but not brute force.

\paragraph{Simulator and real-data grounding.}
\fitStatus\ Validation uses 5-fold cross-validation over participants; a first, contiguous split (participants 1--30 versus 31--40) failed on the congruency effect because the two ID ranges differ, and this change is recorded in the preregistration's deviations log.
Response times follow an ex-Gaussian with participant-level variation; accuracy is logistic; injected effects vary across participants; blinks contaminate frontal theta unless ICA is applied.
In the semi-synthetic track each simulated participant is a real ERP CORE participant, with known effects injected on top of their real responses, theta and blinks.

\fig{fig_validation}{1.0}{Posterior-predictive check: summaries of held-out real ERP CORE participants (points) against simulated participants (histogram). $d$ is the standardised difference.}

\section{Agents, analyst and scoring}
\paragraph{Reference agents.}
All three use the same ideal analyst, so they differ only in experiment choice: \emph{random} (random affordable designs), \emph{scripted} (a textbook counterbalanced $2{\times}2$ with EEG and ICA) and \emph{oracle} (the single design that maximises the ideal analyst's expected accuracy over the hypothesis family, found by simulation).

\paragraph{Ideal analyst.}
A cluster-robust regression per edge, trusting performance effects only from designs not tied to session order and theta effects only from blink-free trials.
Posterior probabilities use the BIC approximation to the Bayes factor, $\log \mathrm{BF}_{10} \approx (z^2 - \ln n)/2$, with $z$ adjusted for the number of participants.

\paragraph{Language-model agents.}
\nModels{} open-weight models (3B to 15B parameters) run locally with schema-constrained JSON actions, a 12-turn limit and fixed sampling temperature, each with variable names shown either as named (``load'', ``tES'') or abstract (``Factor A''), to test reliance on textbook priors.

\paragraph{Metrics.}
The primary metric is the mean edge Brier score, a proper scoring rule that rewards calibrated confidence.
Secondary metrics are exact mechanism recovery, falling for the planted trap, confident errors and whether the agent took the confound-revealing action.
The \emph{ideal-on-own-data} score re-analyses an agent's own datasets with the ideal analyst: the gap from the agent's claims is analysis error, and the gap from the oracle is design error.

\section{Axiom-Bio: real single cells}
\paragraph{Data.} Eleven phosphoproteins and phospholipids measured by flow cytometry in 7{,}466 human CD4$^+$ T cells under nine conditions \citep{sachs2005causal}: CD3/CD28 stimulation, with ICAM-2 co-stimulation, and with five inhibitors (of Akt, PKC, PIP2, MEK and PI3K), plus PKC activation by PMA and PKA activation by $\beta_2$cAMP, both without CD3/CD28. Ground truth is the 20-edge consensus graph distributed with the data (Zenodo record 7681811).
\paragraph{Tasks.} Forty tasks, each with six candidate directed edges: two or three true edges plus reversed, indirect (a path through another molecule) and confounded (a shared regulator, no path) edges. Agents buy real cells (2{,}400 per task, with a per-condition overhead), compare conditions and compute partial correlations.
\paragraph{What it can and cannot measure.} Two real pharmacological subtleties are in the data: U0126, a MEK inhibitor, \emph{raises} measured phospho-Mek while lowering Erk, and the PMA and $\beta_2$cAMP conditions lack the CD3/CD28 context of the baseline. Most consensus edges cannot be identified from the available interventions alone: a data-only reference analyst recovers 22\% of true edges. The track therefore isolates a different question. With real names a model can recall the textbook pathway; with abstract labels it can only read the data. The named-minus-abstract gap (B1) measures recall, and the abstract-label comparison with the data-only analyst (B2) measures inference.

\fig{fig_bio}{0.85}{Axiom-Bio: each model with real molecule names and with abstract labels.}

\section{Preregistered hypotheses}
H1: models score below the scripted baseline.
H2: models fall for planted traps more often than they make the matched error on mirror tasks.
H3: both design and analysis error are non-zero, quantified by the decomposition.
H4: performance improves with parameter count (descriptive, five models).
H5: agents rank the same on synthetic and semi-synthetic tracks (Kendall $\tau \ge 0.6$).
H6: named variables help compared with abstract labels.
Axiom-Bio adds B1 (names help, i.e.\ recall), B2 (abstract-label models versus the data-only analyst) and B3 (claim rates for true, reversed, indirect and confounded edges).
Full text, analysis plan, the stopping rule and a dated deviations log are in \texttt{PREREGISTRATION.md}.

\section{Results}
\resultsStatus

\begin{table}[h]\centering\small
\begin{tabular}{llrrrrr}
\toprule
Agent & Track & $n$ & Brier $\downarrow$ & Exact & Trap & Ideal-on-own \\
\midrule
random & synthetic & 180 & 0.120 & 0.46 & 0.08 & 0.46 \\
random & semisynthetic & 180 & 0.098 & 0.57 & 0.08 & 0.57 \\
scripted & synthetic & 180 & 0.050 & 0.74 & 0.12 & 0.74 \\
scripted & semisynthetic & 180 & 0.049 & 0.76 & 0.15 & 0.76 \\
oracle & synthetic & 180 & 0.041 & 0.82 & 0.10 & 0.82 \\
oracle & semisynthetic & 180 & 0.044 & 0.77 & 0.18 & 0.77 \\
\bottomrule
\end{tabular}

\caption{Named-label results. Brier: mean edge Brier score (lower is better). Exact: fraction of episodes recovering the whole mechanism. Trap: fraction of confound-task episodes claiming the confounded edge. Ideal-on-own: exact recovery when the ideal analyst re-analyses the agent's own data.}
\end{table}

\fig{fig_agents}{0.85}{Primary metric by agent and track, with 95\% bootstrap intervals over tasks.}
\fig{fig_decomposition}{0.85}{Design versus analysis error for each model.}
\fig{fig_traps}{0.85}{How often each agent claimed the planted confounded effect.}
\fig{fig_scale}{0.75}{Scientific judgement against model size.}

\section{Limitations}
The hypothesis family is small and hand-designed; a model could in principle learn its menu, which is why two task types are held out and an abstract-label condition is run.
Effects of load and stimulation are not in ERP CORE; their sizes come from deliberately small priors and are varied rather than fitted.
The ideal analyst knows the practice-curve shape, so the oracle line is an upper reference, not a bound on any possible analyst.
Only local open-weight models up to 15B parameters are evaluated, because the study is run at zero cost.

\section{Reproducibility}
Code, the frozen task manifest and its hash, the fitted parameters, every episode transcript and the results ledger are public.
The figures, this paper's numbers and the project page are regenerated from the ledger on every run.

\bibliographystyle{plainnat}
\bibliography{refs}
\end{document}
